% SpqrBridge slides, randomness alternatives.
% Build: latexmk -lualatex slides.tex
% Shared preamble of the SpqrBridge slide decks.
% Same style and colors as ref-impl-simulation/slides/methodology.
% Each deck loads it with % Shared preamble of the SpqrBridge slide decks.
% Same style and colors as ref-impl-simulation/slides/methodology.
% Each deck loads it with % Shared preamble of the SpqrBridge slide decks.
% Same style and colors as ref-impl-simulation/slides/methodology.
% Each deck loads it with \input{../common/preamble} and builds with: latexmk -lualatex slides.tex
\documentclass[aspectratio=169,11pt]{beamer}

\usepackage{fontspec}
\setsansfont{DejaVu Sans}
\directlua{
  luaotfload.add_fallback("leanfallback", {"NotoSansMath:mode=harf;", "STIX Two Math:mode=harf;"})
}
\setmonofont{DejaVu Sans Mono}[Scale=0.9, RawFeature={fallback=leanfallback}]

\usepackage{tikz}
\usepackage{array}
\usetikzlibrary{arrows.meta, positioning, calc, fit, backgrounds, shapes.misc, decorations.pathreplacing}

% Colors. Spec side, impl side, proof, trusted, gap.
\definecolor{refc}{HTML}{2B6CB0}
\definecolor{realc}{HTML}{C05621}
\definecolor{proofc}{HTML}{2F855A}
\definecolor{trustc}{HTML}{718096}
\definecolor{gapc}{HTML}{C53030}
\definecolor{ink}{HTML}{1A202C}
\definecolor{paper}{HTML}{FFFFFF}

\setbeamercolor{normal text}{fg=ink, bg=paper}
\setbeamercolor{frametitle}{fg=ink}
\setbeamercolor{title}{fg=ink}
\setbeamercolor{subtitle}{fg=trustc}
\setbeamercolor{structure}{fg=refc}
\setbeamerfont{frametitle}{size=\Large, series=\bfseries}
\setbeamerfont{subtitle}{size=\footnotesize}
\setbeamertemplate{navigation symbols}{}
\setbeamertemplate{footline}{%
  \hfill\textcolor{trustc}{\scriptsize\insertframenumber/\inserttotalframenumber}\hspace{1.2em}\vspace{0.8em}}

\tikzset{
  >={Stealth[length=2.4mm]},
  box/.style={draw, rounded corners=3pt, thick, align=center, inner sep=5pt, font=\small},
  spec/.style={box, draw=refc, fill=refc!8},
  impl/.style={box, draw=realc, fill=realc!8},
  proof/.style={box, draw=proofc, fill=proofc!10},
  trust/.style={box, draw=trustc, fill=trustc!10},
  gap/.style={box, draw=gapc, fill=gapc!6, dashed},
  tool/.style={box, draw=ink!60, fill=ink!4},
  doc/.style={box, draw=ink!60, fill=paper, double, double distance=1pt},
  flow/.style={->, thick, ink!70},
  proved/.style={<->, very thick, proofc},
  rel/.style={-, very thick, proofc, dashed},
  label/.style={font=\scriptsize, align=center},
}

\newcommand{\code}[1]{\texttt{#1}}
\newcommand{\specside}[1]{\textcolor{refc}{#1}}
\newcommand{\implside}[1]{\textcolor{realc}{#1}}

\author{SpqrBridge}
\date{November 2026}
 and builds with: latexmk -lualatex slides.tex
\documentclass[aspectratio=169,11pt]{beamer}

\usepackage{fontspec}
\setsansfont{DejaVu Sans}
\directlua{
  luaotfload.add_fallback("leanfallback", {"NotoSansMath:mode=harf;", "STIX Two Math:mode=harf;"})
}
\setmonofont{DejaVu Sans Mono}[Scale=0.9, RawFeature={fallback=leanfallback}]

\usepackage{tikz}
\usepackage{array}
\usetikzlibrary{arrows.meta, positioning, calc, fit, backgrounds, shapes.misc, decorations.pathreplacing}

% Colors. Spec side, impl side, proof, trusted, gap.
\definecolor{refc}{HTML}{2B6CB0}
\definecolor{realc}{HTML}{C05621}
\definecolor{proofc}{HTML}{2F855A}
\definecolor{trustc}{HTML}{718096}
\definecolor{gapc}{HTML}{C53030}
\definecolor{ink}{HTML}{1A202C}
\definecolor{paper}{HTML}{FFFFFF}

\setbeamercolor{normal text}{fg=ink, bg=paper}
\setbeamercolor{frametitle}{fg=ink}
\setbeamercolor{title}{fg=ink}
\setbeamercolor{subtitle}{fg=trustc}
\setbeamercolor{structure}{fg=refc}
\setbeamerfont{frametitle}{size=\Large, series=\bfseries}
\setbeamerfont{subtitle}{size=\footnotesize}
\setbeamertemplate{navigation symbols}{}
\setbeamertemplate{footline}{%
  \hfill\textcolor{trustc}{\scriptsize\insertframenumber/\inserttotalframenumber}\hspace{1.2em}\vspace{0.8em}}

\tikzset{
  >={Stealth[length=2.4mm]},
  box/.style={draw, rounded corners=3pt, thick, align=center, inner sep=5pt, font=\small},
  spec/.style={box, draw=refc, fill=refc!8},
  impl/.style={box, draw=realc, fill=realc!8},
  proof/.style={box, draw=proofc, fill=proofc!10},
  trust/.style={box, draw=trustc, fill=trustc!10},
  gap/.style={box, draw=gapc, fill=gapc!6, dashed},
  tool/.style={box, draw=ink!60, fill=ink!4},
  doc/.style={box, draw=ink!60, fill=paper, double, double distance=1pt},
  flow/.style={->, thick, ink!70},
  proved/.style={<->, very thick, proofc},
  rel/.style={-, very thick, proofc, dashed},
  label/.style={font=\scriptsize, align=center},
}

\newcommand{\code}[1]{\texttt{#1}}
\newcommand{\specside}[1]{\textcolor{refc}{#1}}
\newcommand{\implside}[1]{\textcolor{realc}{#1}}

\author{SpqrBridge}
\date{November 2026}
 and builds with: latexmk -lualatex slides.tex
\documentclass[aspectratio=169,11pt]{beamer}

\usepackage{fontspec}
\setsansfont{DejaVu Sans}
\directlua{
  luaotfload.add_fallback("leanfallback", {"NotoSansMath:mode=harf;", "STIX Two Math:mode=harf;"})
}
\setmonofont{DejaVu Sans Mono}[Scale=0.9, RawFeature={fallback=leanfallback}]

\usepackage{tikz}
\usepackage{array}
\usetikzlibrary{arrows.meta, positioning, calc, fit, backgrounds, shapes.misc, decorations.pathreplacing}

% Colors. Spec side, impl side, proof, trusted, gap.
\definecolor{refc}{HTML}{2B6CB0}
\definecolor{realc}{HTML}{C05621}
\definecolor{proofc}{HTML}{2F855A}
\definecolor{trustc}{HTML}{718096}
\definecolor{gapc}{HTML}{C53030}
\definecolor{ink}{HTML}{1A202C}
\definecolor{paper}{HTML}{FFFFFF}

\setbeamercolor{normal text}{fg=ink, bg=paper}
\setbeamercolor{frametitle}{fg=ink}
\setbeamercolor{title}{fg=ink}
\setbeamercolor{subtitle}{fg=trustc}
\setbeamercolor{structure}{fg=refc}
\setbeamerfont{frametitle}{size=\Large, series=\bfseries}
\setbeamerfont{subtitle}{size=\footnotesize}
\setbeamertemplate{navigation symbols}{}
\setbeamertemplate{footline}{%
  \hfill\textcolor{trustc}{\scriptsize\insertframenumber/\inserttotalframenumber}\hspace{1.2em}\vspace{0.8em}}

\tikzset{
  >={Stealth[length=2.4mm]},
  box/.style={draw, rounded corners=3pt, thick, align=center, inner sep=5pt, font=\small},
  spec/.style={box, draw=refc, fill=refc!8},
  impl/.style={box, draw=realc, fill=realc!8},
  proof/.style={box, draw=proofc, fill=proofc!10},
  trust/.style={box, draw=trustc, fill=trustc!10},
  gap/.style={box, draw=gapc, fill=gapc!6, dashed},
  tool/.style={box, draw=ink!60, fill=ink!4},
  doc/.style={box, draw=ink!60, fill=paper, double, double distance=1pt},
  flow/.style={->, thick, ink!70},
  proved/.style={<->, very thick, proofc},
  rel/.style={-, very thick, proofc, dashed},
  label/.style={font=\scriptsize, align=center},
}

\newcommand{\code}[1]{\texttt{#1}}
\newcommand{\specside}[1]{\textcolor{refc}{#1}}
\newcommand{\implside}[1]{\textcolor{realc}{#1}}

\author{SpqrBridge}
\date{November 2026}


\title{Random bytes in the bridge: four approaches}
\subtitle{SpqrBridge slides, randomness alternatives}

\tikzset{
  cell/.style={draw=realc, fill=realc!8, minimum width=0.42cm, minimum height=0.42cm,
    inner sep=0pt, font=\tiny},
  readcell/.style={cell, fill=realc!35},
  box/.append style={font=\footnotesize},
  line/.style={font=\scriptsize\ttfamily, anchor=west, inner sep=2pt},
}
\newcommand{\source}[1]{{\footnotesize\textcolor{trustc}{#1}}\par\vspace{0.2em}}
\newcommand{\slidenote}[1]{\par\vspace{0.6em}\centering{\footnotesize #1}\par}
\newcommand{\pro}{\textcolor{proofc}{\textbf{+}}}
\newcommand{\con}{\textcolor{gapc}{\textbf{--}}}
\newcommand{\yes}{\textcolor{proofc}{\textbf{yes}}}
\newcommand{\no}{\textcolor{gapc}{\textbf{no}}}
% Pros and cons in two columns: \proscons{pros}{cons}
\newcommand{\proscons}[2]{%
  \raggedright
  \begin{columns}[T]
    \column{0.42\textwidth}\footnotesize
    \textbf{\textcolor{proofc}{For}}\\[0.3em] #1
    \column{0.42\textwidth}\footnotesize
    \textbf{\textcolor{gapc}{Against}}\\[0.3em] #2
  \end{columns}}

\begin{document}

% ---------------------------------------------------------------------------
\begin{frame}
  \titlepage
  \centering
  {\footnotesize\textcolor{trustc}{Aeneas code is deterministic. VCVio proofs need
  distributions. Where does randomness enter?}}
\end{frame}

% ---------------------------------------------------------------------------
\begin{frame}[t]{The running example: \code{mask}}
  \vfill
  \centering
  \begin{tikzpicture}[
      rs/.style={font=\scriptsize\ttfamily, anchor=west, inner sep=1pt},
      mid/.style={font=\scriptsize, text=ink!70, anchor=center}]
    % highlight of the one row that differs
    \fill[gapc!10, rounded corners=2pt] (-6.6, -1.28) rectangle (6.6, -0.72);
    % Rust
    \node[font=\small\bfseries, text=realc, anchor=west] at (-6.4, 0.75) {Rust \code{mask}};
    \node[rs] at (-6.4, 0.0)  {fn mask(rng: \&mut R, v: u8)};
    \node[rs] at (-6.4, -0.5) {\ \ let mut pad = [0u8; 1];};
    \node[rs] at (-6.4, -1.0) {\ \ rng.fill\_bytes(\&mut pad);};
    \node[rs] at (-6.4, -1.5) {\ \ (pad[0], v \textasciicircum{} pad[0])};
    \begin{scope}[on background layer]
      \draw[realc, thick, rounded corners=3pt, fill=realc!5] (-6.55, 0.35) rectangle (-1.75, -1.85);
    \end{scope}
    % VCVio
    \node[font=\small\bfseries, text=refc, anchor=west] at (1.95, 0.75) {VCVio \code{Crypto.mask}};
    \node[rs] at (1.95, 0.0)  {def mask (v : U8) :=};
    \node[rs] at (1.95, -1.0) {\ \ let pad ← \$ᵗ U8};
    \node[rs] at (1.95, -1.5) {\ \ return (pad, v ⊕ pad)};
    \begin{scope}[on background layer]
      \draw[refc, thick, rounded corners=3pt, fill=refc!5] (1.8, 0.35) rectangle (6.55, -1.85);
    \end{scope}
    % middle column: how the rows relate
    \node[mid] at (0, 0.0)  {input \code{v}};
    \node[mid, text=gapc] at (0, -1.0) {\textbf{get the byte \code{pad}}};
    \node[mid] at (0, -1.5) {same output};
    % the difference, spelled out
    \node[font=\scriptsize, text=realc, align=center] at (-4.15, -2.5)
      {\code{pad} = what the generator writes\\one fixed answer};
    \node[font=\scriptsize, text=refc, align=center] at (4.15, -2.5)
      {\code{pad} = uniform random byte\\a distribution over 256 answers};
  \end{tikzpicture}

  \slidenote{Same computation. Only the line that gets \code{pad} differs.\\[0.3em]
    \code{Crypto.mask} is what \code{mask} computes when the generator is perfect.}
  \vfill
\end{frame}

% ---------------------------------------------------------------------------
\begin{frame}[t]{Where the bytes really come from}
  \vfill
  \centering
  \begin{tikzpicture}[every node/.append style={font=\footnotesize}]
    \node[tool, text width=1.2cm, minimum height=1.3cm] (app) at (0, 0) {app};
    \node[impl, text width=2.6cm, minimum height=1.3cm] (entry) at (2.9, 0) {libsignal\\creates \code{OsRng}};
    \node[impl, text width=2.8cm, minimum height=1.3cm] (send) at (6.9, 0) {\code{spqr::send<R>}\\calls \code{fill\_bytes}};
    \node[trust, text width=1.4cm, minimum height=1.3cm] (os) at (11.9, 0) {kernel};
    \draw[flow] (app) -- (entry);
    \draw[flow] (entry) -- node[above, label] {\code{OsRng}} (send);
    \draw[flow, trustc] ($(send.east)+(0,0.25)$) -- node[above, label, text=trustc] {system call} ($(os.west)+(0,0.25)$);
    \draw[flow, trustc] ($(os.west)+(0,-0.25)$) -- node[below, label, text=trustc] {random bytes} ($(send.east)+(0,-0.25)$);
    \draw[decorate, decoration={brace, amplitude=5pt, mirror}, realc, thick]
      (entry.west |- 0,-0.95) -- (send.east |- 0,-0.95)
      node[midway, below=6pt, label, text=realc] {Aeneas sees this};
    \draw[decorate, decoration={brace, amplitude=5pt, mirror}, trustc, thick]
      (os.west |- 0,-0.95) -- (os.east |- 0,-0.95)
      node[midway, below=6pt, label, text=trustc] {outside Lean};
  \end{tikzpicture}
  \vfill
  \slidenote{Inside \code{send}, each \code{fill\_bytes} on \code{OsRng} is a system call.\\[0.3em]
    \code{OsRng} holds no state. Every approach trusts:
    \textcolor{trustc}{``the kernel's bytes are uniform''}.}
  \vfill
\end{frame}

% ---------------------------------------------------------------------------
\begin{frame}[t]{The problem}
  \vfill
  \centering
  \begin{tikzpicture}[every node/.append style={font=\footnotesize}]
    \node[impl, text width=4.2cm, minimum height=2.0cm] (l) at (0, 0) {%
      \textbf{Aeneas}\\[4pt]
      \code{mask inst rng v : Result \_}\\[4pt]
      same \code{rng} in, same bytes out};
    \node[spec, text width=4.2cm, minimum height=2.0cm] (r) at (9.6, 0) {%
      \textbf{VCVio}\\[4pt]
      \code{Crypto.mask v : ProbComp \_}\\[4pt]
      a distribution};
    \node[gap, text width=2.0cm] (g) at (4.8, 0) {\textbf{?}};
    \draw[gapc, dashed, thick] (l) -- (g) -- (r);
  \end{tikzpicture}

  \slidenote{Goal: \emph{if} the generator gives uniform bytes, \emph{then} Rust \code{mask} $=$
    \code{Crypto.mask}.\\[0.3em] Next, we go through several approaches that differ in how they state that ``if''.}
  \vfill
\end{frame}

% ---------------------------------------------------------------------------
\begin{frame}[t]{How we compare them}
  \vfill
  \centering
  \begin{tikzpicture}
    \node[box, font=\small, text width=6.6cm, minimum height=1.2cm] at (-3.6, 0) {\textbf{Q1.} Is the ``if'' a Lean statement,\\in one theorem?};
    \node[box, font=\small, text width=6.6cm, minimum height=1.2cm] at (3.6, 0) {\textbf{Q2.} How much extra work\\per function?};
  \end{tikzpicture}
  \vfill
  {\color{trustc!50}\rule{0.85\textwidth}{0.4pt}}
  \vfill
  {\footnotesize\textcolor{trustc}{How to read the diagrams}}\par\vspace{0.6em}
  \begin{tikzpicture}[every node/.append style={font=\footnotesize}]
    \node[impl, minimum width=2.6cm] at (0, 0) {extracted code};
    \node[spec, minimum width=2.6cm] at (3.2, 0) {VCVio};
    \node[proof, minimum width=2.6cm] at (6.4, 0) {written by us};
    \node[trust, minimum width=2.6cm] at (0, -1.0) {assumption};
    \node[gap, minimum width=2.6cm] at (3.2, -1.0) {missing};
    \draw[proved] (5.3, -1.0) -- (6.3, -1.0) node[right, label, font=\footnotesize, text=proofc] {proved equal};
  \end{tikzpicture}
  \vfill
\end{frame}

% ===========================================================================
\begin{frame}[t]{1. Random tape}
  \source{spqr-sm-bridge before \#10. Replaced.}
  \vfill
  \centering
  \begin{tikzpicture}[every node/.append style={font=\footnotesize}]
    \node[spec, text width=2.0cm] (samp) at (0, 0) {sample\\64 bytes};
    \foreach \i in {0,...,7} {
      \ifnum\i<1
        \node[readcell] (c\i) at (2.9+0.42*\i, 0) {};
      \else
        \node[cell] (c\i) at (2.9+0.42*\i, 0) {};
      \fi
    }
    \node at (6.5, 0) {$\cdots$};
    \node[cell] (clast) at (7.0, 0) {};
    \draw[flow] (samp) -- (c0.west);
    \node[label, text=realc, anchor=south] at (c0.north) {read};
    \draw[decorate, decoration={brace, amplitude=4pt, mirror}, gapc, thick]
      ($(c1.south west)+(0,-0.1)$) -- ($(clast.south east)+(0,-0.1)$)
      node[midway, below=4pt, label, text=gapc] {never read};
    \node[impl, text width=2.4cm] (run) at (10.0, 0) {\code{mask} on the tape};
    \draw[flow] (clast.east) -- (run);
    \node[spec, text width=2.4cm] (cr) at (10.0, -1.8) {\code{Crypto.mask}};
    \draw[proved] (run) -- (cr);
  \end{tikzpicture}

  \slidenote{Roll all the dice first. The code reads them.}
  \vfill
  \proscons{%
    \pro\ No change to Aeneas\\[0.2em]
    \pro\ Simple}{%
    \con\ Q1: no. The ``if'' is prose\\[0.2em]
    \con\ Q2: a tape length per function\\[0.2em]
    \con\ Q2: lemmas for unread bytes}
  \vfill
\end{frame}

% ===========================================================================
\begin{frame}[t]{2. Contract on the generator state}
  \source{bridging-examples \code{main}. Removed in \#1.}
  \vfill
  \centering
  \begin{tikzpicture}[every node/.append style={font=\footnotesize}]
    \node[trust, minimum width=2.0cm] (s2) at (0, 0) {state, 2 left};
    \node[trust, minimum width=2.0cm] (s1) at (5.0, 0) {state, 1 left};
    \node[trust, minimum width=2.0cm] (s0) at (10.0, 0) {state, 0 left};
    \node[spec] (b1) at (2.5, -1.2) {uniform byte};
    \node[spec] (b2) at (7.5, -1.2) {uniform byte};
    \draw[flow, realc] (s2) -- node[above, label, text=realc] {draw} (s1);
    \draw[flow, realc] (s1) -- node[above, label, text=realc] {draw} (s0);
    \draw[flow, refc] ($(s2)!0.5!(s1)$) -- (b1);
    \draw[flow, refc] ($(s1)!0.5!(s0)$) -- (b2);
  \end{tikzpicture}

  \slidenote{Assume what the state looks like before each draw (\code{RngContract}).}
  \vfill
  \proscons{%
    \pro\ No change to Aeneas\\[0.2em]
    \pro\ Q1: yes}{%
    \con\ Q2: a draw count per function\\[0.2em]
    \con\ Randomness only from the state.\\\ \ \ \code{OsRng} keeps it in the OS}
  \vfill
\end{frame}

% ===========================================================================
\begin{frame}[t]{3. Oracle program\footnotemark[1] + \code{IdealRng}}
  \source{\scriptsize bridging-examples\#3, spqr-sm-bridge\#10}
  \footnotetext[1]{\tiny\href{https://github.com/Beneficial-AI-Foundation/toy-bridge}{toy-bridge}: the oracle idea behind \#10 and \#3, with the function copied by hand and without \code{IdealRng}.}
  \vfill
  \centering
  \begin{tikzpicture}[every node/.append style={font=\footnotesize}]
    \node[proof, text width=7.2cm] (p) at (0, 0) {\code{maskProg v = b <- query ; maskRest v b}};
    \node[impl, text width=2.8cm] (l) at (-4.6, -1.2) {Rust \code{mask}};
    \node[spec, text width=2.8cm] (r) at (4.6, -1.2) {\code{Crypto.mask}};
    \draw[proved] (p.west) -| (l.north);
    \draw[proved] (p.east) -| (r.north);
    \draw[proved, line width=2pt] (l.south) -- ++(0,-0.45) -| (r.south);
    \node[label, text=proofc, above] at (0, -2.05) {\textbf{one theorem}};
    \node[trust, text width=8.4cm] (h) at (0, -2.6) {\code{IdealRng}: the generator answers like uniform bytes};
    \draw[flow, trustc] (h.north) -- (0, -2.05);
  \end{tikzpicture}

  \slidenote{\code{maskRest v b}: the unchanged extracted \code{mask}, with a generator that returns \code{b}.\\
    So the bytes are an argument (\code{sendRest} in \#10).}
  \vspace{-0.4em}
  \proscons{%
    \scriptsize
    \pro\ No change to Aeneas\\[0.1em]
    \pro\ Q1: yes\\[0.1em]
    \pro\ \code{IdealRng} is not empty:\\\ \ \ a uniform tape meets it}{%
    \scriptsize
    \con\ Q2: a program per function\\[0.1em]
    \con\ A budget, one state per program\\[0.1em]
    \con\ Randomness must come from the state.\\\ \ \ \code{OsRng} keeps it in the OS}
  \vfill
\end{frame}

% ===========================================================================
\begin{frame}[t]{4. Sampling effect in \code{Result}}
  \source{Sergiu, proposal}
  \vfill
  \centering
  \begin{tikzpicture}[every node/.append style={font=\footnotesize}]
    \node[gap, text width=6.4cm] (patch) at (0, 0) {patch Aeneas: a \code{Result} can draw a uniform value};
    \node[impl, text width=2.8cm] (l) at (-4.6, -1.3) {Rust \code{mask}};
    \node[spec, text width=2.8cm] (r) at (4.6, -1.3) {\code{Crypto.mask}};
    \draw[proved, line width=2pt] (l) -- node[pos=0.22, below, label, text=proofc] {\textbf{one theorem}} (r);
    \node[trust, text width=5.0cm] (h) at (0, -2.5) {H(G): \code{fill\_bytes} draws uniform bytes};
    \draw[flow, trustc] (h.north) -- (0, -1.45);
    \draw[flow, gapc, dashed] (patch.south) -- (0, -1.15);
  \end{tikzpicture}

  \slidenote{Roll the dice inside the code, at the call.}
  \vfill
  \proscons{%
    \pro\ Q1: yes\\[0.2em]
    \pro\ Q2: no extra program\\[0.2em]
    \pro\ Fits a stateless \code{OsRng}}{%
    \con\ Changes \code{Result} for all users\\[0.2em]
    \con\ Functional proofs stop at a draw\\[0.2em]
    \con\ First patch withdrawn, new one not seen}
  \vfill
\end{frame}

% ===========================================================================
\begin{frame}[t]{3 and 4: same trust, different shape}
  \vfill
  \centering
  \begin{tikzpicture}[every node/.append style={font=\footnotesize}]
    \node[font=\small\bfseries, text=proofc] at (-3.4, 0.6) {3 \code{IdealRng}};
    \node[trust, text width=5.0cm] at (-3.4, -0.3) {all answers of a program\\look uniform};
    \node[font=\small\bfseries, text=gapc] at (3.4, 0.6) {4 H(G)};
    \node[trust, text width=5.0cm] at (3.4, -0.3) {each call\\draws uniform bytes};
    \node[label, text width=5.0cm] at (-3.4, -1.6) {global, with a budget\\Aeneas unchanged};
    \node[label, text width=5.0cm] at (3.4, -1.6) {local, no budget\\Aeneas patched};
  \end{tikzpicture}

  \slidenote{In both, ``\code{OsRng} meets the assumption'' stays outside Lean.}
  \vfill
\end{frame}

% ---------------------------------------------------------------------------
\begin{frame}[t]{Side by side}
  \vfill
  \centering
  \footnotesize
  \renewcommand{\arraystretch}{1.6}
  \setlength{\tabcolsep}{5pt}
  \begin{tabular}{l c c c}
    & \textbf{Aeneas unchanged} & \textbf{Q1: ``if'' in Lean} & \textbf{Q2: extra work}\\
    \hline
    1 (Tape)          & \yes & \no  & tape length \\
    2 (Contract)      & \yes & \yes & draw counts \\
    3 (Oracle)        & \yes & \yes & program per function \\
    4 (Effect)        & \no  & \yes & none \\
    \hline
  \end{tabular}
  \vfill
\end{frame}

\end{document}
